\section{Recall analysis}
\label{sec:recall}
A search that returns 80 candidates from a database of 256 reduces the remaining scoring
work to those candidates. Recall depends on how many of the nearest documents are among
those 80. This section derives both quantities under a small random-bit model, then extends
the analysis to query weights and clusters.

\Needspace{12\baselineskip}
\subsection{Recall definition}
Let $\mathcal E(q)$ be the eligible document set after the intended filter, across the whole
database. Let $\mathcal G_K(q)$ be its exact top $K$ under the chosen reference score, and let
$\mathcal R_K(q)$ be the returned results. Then
\begin{equation}
\operatorname{Recall@}K(q)=
\frac{|\mathcal G_K(q)\cap\mathcal R_K(q)|}{K}.
\label{eq:recall-exact}
\end{equation}
Both sets use the same tie rule, and mean recall averages this quantity over queries.
Repeated timing measurements of one query do not provide independent quality observations.

If a candidate set $\mathcal D(q)$ is reranked with the \emph{same exact reference score}, every
true top-$K$ document in that set survives into the final top $K$. Therefore
\begin{equation}
\operatorname{Recall@}K(q)=\frac{|\mathcal G_K(q)\cap\mathcal D(q)|}{K}.
\label{eq:candidate-recall}
\end{equation}
A different reranking score breaks that identity. In particular, exact binary candidates may
still miss the true float-vector top $K$ or the documents relevant to a human question.

\fig{recall-flow}{Follow the same true neighbors through every stage. The denominator stays K throughout the recall calculation.}{.76}

\Needspace{12\baselineskip}
\subsection{A small random-bit model}
Suppose a document has $d$ bits. We sample $N$ documents independently, and each bit is
independently zero or one with probability $1/2$. Duplicate codes are allowed: two different
documents can have the same bits. Their row IDs resolve score ties, and those IDs are assigned
independently of the sampled codes.

Initially, every query coordinate has the same magnitude. A document's score therefore depends
only on how many bits disagree with the query. We can write the query's desired pattern as all
zeros without changing any probabilities. The number of mismatches is called the
\emph{Hamming distance}. A smaller distance means a better score.

The search checks $h$ fixed coordinates and accepts a document if at most $t$ of them disagree.
It then scores every accepted document using all $d$ bits. For example, with $h=4$ and $t=1$,
it opens the keys $0000$, $0001$, $0010$, $0100$, and $1000$. This is a precise candidate rule.
It does not stop halfway through a key's documents when a quota is reached.

These assumptions define an illustrative model. Equal weights permit direct counting,
and independent bits make every pattern equally likely. Later subsections relax these assumptions.
There are no clusters or eligibility filters in this first calculation.

\Needspace{12\baselineskip}
\subsection{Nearest-document distance}
Let $Z$ be one random document's distance from the query. To have exactly $m$ errors, choose
which $m$ of the $d$ positions are wrong. There are $\binom{d}{m}$ choices, out of $2^d$
equally likely codes. Thus
\begin{equation}
\Pr(Z=m)=\frac{\binom{d}{m}}{2^d},\qquad
F_d(m)=\Pr(Z\le m)=\frac{1}{2^d}\sum_{j=0}^{m}\binom{d}{j}.
\label{eq:one-document-distance}
\end{equation}
This is the binomial distribution with parameters $d$ and $1/2$. The function $F_d$ sums the
probabilities of all distances through $m$. This accumulated probability defines a cumulative
distribution function, or CDF. Define $F_d(-1)=0$.

Let $M$ be the distance of the nearest document among all $N$. If one document has
probability $1-F_d(m)$ of being farther than $m$, independence gives probability
$[1-F_d(m)]^N$ that \emph{all} documents are farther away. Subtracting from one gives
\begin{equation}
\Pr(M\le m)=1-[1-F_d(m)]^N.
\label{eq:nearest-distance-cdf}
\end{equation}
The probability of exactly distance $m$ follows by subtracting the probability of a smaller distance:
\begin{equation}
\Pr(M=m)=[1-F_d(m-1)]^N-[1-F_d(m)]^N.
\label{eq:nearest-distance-mass}
\end{equation}
A random document usually has many mismatched bits, whereas the nearest of hundreds of
documents has unusually few. Recall therefore depends on the distribution of the nearest
distance, rather than the distance of an arbitrary document.

\fig{nearest-distance}{With 256 random 16-bit documents, the nearest distance is usually two or three. Assuming it is always two changes the recall calculation.}{.88}

\Needspace{12\baselineskip}
\subsection{Key acceptance probability}
Suppose the nearest document has exactly $m$ wrong bits. Under the uniform model, every set of
$m$ error positions is equally likely. Let $A$ count how many of those errors fall in the $h$
lookup positions. To have exactly $a$ lookup errors, choose $a$ positions from the lookup bits
and $m-a$ from the other bits. Divide by all possible placements of the $m$ errors:
\begin{equation}
\Pr(A=a\mid M=m)=
\frac{\binom{h}{a}\binom{d-h}{m-a}}{\binom{d}{m}}.
\label{eq:prefix-error-placements}
\end{equation}
This is a hypergeometric distribution. It counts positions drawn from a fixed number of errors.
The error count is already fixed here, so treating these positions as independent trials
would change the model. A binomial coefficient is zero when its lower argument is negative or
larger than its upper argument.

The search accepts at most $t$ lookup errors. Summing the allowed placements gives:
\begin{equation}
g_{h,t}(m)=\Pr(A\le t\mid M=m)
=\frac{\displaystyle\sum_{a=0}^{\min(h,t)}
          \binom{h}{a}\binom{d-h}{m-a}}{\binom{d}{m}}.
\label{eq:prefix-survival}
\end{equation}
Even after we know the nearest distance, every placement of its errors is equally likely because the score
uses only the total number of errors and the ID tie rule does not favor a particular pattern.
This symmetry depends on the equal-weight assumption.

\begin{examplebox}{Two wrong bits in a 16-bit document}
Take $d=16$, $h=4$, and suppose $m=2$. There are $\binom{16}{2}=120$ ways to place the errors.
If the first four bits must match perfectly, both errors must be in the other twelve bits:
$\binom{12}{2}=66$ placements survive. Recall, given two total errors, is $66/120=55\%$.

Allow one error in the first four bits. Another $4\times12=48$ placements now survive, giving
$(66+48)/120=95\%$. This 95\% result is conditional on \emph{two total errors}. It does not yet account
for other possible distances of the nearest document.
\end{examplebox}

\Needspace{12\baselineskip}
\subsection{Expected top-one recall}
For top one, recall is one if the true nearest document is accepted and zero otherwise.
Its expected value is therefore its probability of being accepted. For each possible nearest
distance, multiply the probability of that distance by the probability of passing the prefix
check. Add the results:
\begin{equation}
\boxed{\E[\operatorname{Recall@1}]
=\sum_{m=0}^{d}\Pr(M=m)\,g_{h,t}(m).}
\label{eq:finite-top-one}
\end{equation}
For $N=256$, $d=16$, and $h=4$, this produces the table below. The candidate counts are derived
in the next subsection.

\begin{center}
\begin{tabular}{@{}rrrr@{}}
\toprule
Allowed errors $t$ & Keys opened & Expected documents & Expected Recall@1\\\midrule
0 & 1 & 16 & 46.40\%\\
1 & 5 & 80 & 89.60\%\\
2 & 11 & 176 & 99.45\%\\
4 & 16 & 256 & 100.00\%\\\bottomrule
\end{tabular}
\end{center}
For one allowed lookup error, the result is 89.60\%, rather than the conditional 95\% above.
In this example, the nearest distance is two about 35.05\% of the time and three about
52.05\% of the time. A three-error neighbor is more likely to put too many errors in the
lookup bits. Equation~\eqref{eq:finite-top-one} includes both cases, along with every other
possible distance.

\par\noindent\begin{minipage}{\linewidth}
\begin{lstlisting}[caption={Calculate expected top-one recall under the random-bit assumptions.}]
ExpectedRecall(N, d, h, allowed_errors):
    recall = 0
    for m in 0 .. d:
        nearest_probability = (1 - F_d(m-1))^N - (1 - F_d(m))^N
        passing_placements = 0
        for a in 0 .. min(h, allowed_errors):
            passing_placements += choose(h, a) * choose(d-h, m-a)
        survival = passing_placements / choose(d, m)
        recall += nearest_probability * survival
    return recall
\end{lstlisting}
\end{minipage}\par

\Needspace{12\baselineskip}
\subsection{Candidate and lookup counts}
Opening all $h$-bit keys through radius $t$ requires
\begin{equation}
H(h,t)=\sum_{a=0}^{t}\binom{h}{a},\qquad 0\le t\le h.
\label{eq:teaching-keys}
\end{equation}
Each random document has probability $p_{\rm pass}=H(h,t)/2^h$ of belonging to one of those keys.
If $F$ counts accepted documents, then
\begin{equation}
F\sim\operatorname{Binomial}(N,p_{\rm pass}),\qquad
\E[F]=Np_{\rm pass},\qquad
\Var(F)=Np_{\rm pass}(1-p_{\rm pass}).
\label{eq:teaching-candidates}
\end{equation}
For $t=1$, five of sixteen keys are opened. The expected candidate count is
$256\times5/16=80$, with standard deviation $\sqrt{55}\approx7.42$ documents. A particular
corpus need not return exactly 80. The 31.25\% candidate fraction also differs from the 89.60\%
recall: nearest documents are more likely than arbitrary documents to pass this filter.

For C, $H$ counts lookup attempts, including attempts that return nothing. The documents in
those keys supply the posting and scoring work. A short-key query that scores all accepted
rows might have expected time
\[
\E[T_C]=T_{\rm prepare}+Hc_{\rm lookup}
       +\E[F](c_{\rm post}+c_{\rm score})+\E[T_{\rm select}],
\]
with enumeration and routing added when the implementation performs them. Selection time can
depend nonlinearly on the returned rows; substituting $\E[F]$ into every time term is another
approximation. Use the full cost model in Equation~\eqref{eq:final-C} for the actual layout.

Storing every candidate ID at once takes $F b_{\rm id}$ bytes. Streaming postings through a
bounded top-$K$ heap can avoid that particular full candidate array. Neither choice removes
the index itself, and an expected candidate count is not a peak-RAM guarantee.

\Needspace{12\baselineskip}
\subsection{Expected top-K recall}
For top $K$, the $K$ selected rows can lie at different distances: some may have zero errors,
some one, and others more. The calculation therefore requires the expected number of selected
rows at each distance.
Assume $1\le K\le N$.

Let $X_m$ count corpus documents whose distance is at most $m$. Since every document independently
has probability $F_d(m)$ of meeting that condition,
\[
X_m\sim\operatorname{Binomial}(N,F_d(m)).
\]
Exactly $\min(K,X_m)$ of the true top-$K$ rows lie within that radius. There cannot be more
than $K$, and when fewer than $K$ documents qualify, every qualifying document belongs to the
top $K$. Define
\begin{equation}
G_K(p)=\E[\min(K,X)],\quad X\sim\operatorname{Binomial}(N,p).
\label{eq:clipped-binomial}
\end{equation}
Then the expected number of top-$K$ rows at exactly distance $m$ is
\[
w_m=G_K(F_d(m))-G_K(F_d(m-1)).
\]
The same placement calculation applies to each row at that distance. Therefore
\begin{equation}
\boxed{\E[\operatorname{Recall@}K]
=\frac{1}{K}\sum_{m=0}^{d}w_m g_{h,t}(m).}
\label{eq:finite-top-k}
\end{equation}
This includes tied distances and duplicate codes under the stated ID rule. For $K=1$, $G_1(p)$ is
$1-(1-p)^N$, so it reduces to Equation~\eqref{eq:finite-top-one}. For $K=N$, every document is
in the reference set and recall becomes the candidate fraction. Opening every key gives
recall one for every $K$.

The appendix gives a numerical method for evaluating $G_K$. The formula accounts for the
varying $K$-th distance without replacing it with a fixed estimated radius. It assumes a fixed
prefix acceptance rule. Stopping after the first $R$ postings is a
different rule whose boundary depends on the sampled corpus.

\Needspace{12\baselineskip}
\subsection{Routing and local recall}
Let $\mathcal U(q)$ be the eligible pool in the probed clusters. Define
\[
r_{\rm route}(q)=\frac{|\mathcal G_K(q)\cap\mathcal U(q)|}{K},\qquad
r_{\rm local}(q)=\frac{|\mathcal G_K(q)\cap\mathcal D(q)|}
 {|\mathcal G_K(q)\cap\mathcal U(q)|}.
\]
When the denominator is positive and $\mathcal D\subseteq\mathcal U$,
\begin{equation}
\operatorname{Recall@}K(q)=r_{\rm route}(q)r_{\rm local}(q).
\label{eq:recall-stages}
\end{equation}
If no true top-$K$ result enters the selected pool, total recall is zero. This multiplication is
exact when both factors refer to the \emph{same true global neighbors}. However, multiplying
an average routing recall by an average local recall is generally incorrect:
\[
\E[r_{\rm route}r_{\rm local}]
=\E[r_{\rm route}]\E[r_{\rm local}]+\Cov(r_{\rm route},r_{\rm local}).
\]
The covariance term records whether the two recall values tend to rise or fall together across queries.
Recall against a cluster's local top $K$ can also differ from recall for the particular
global neighbors routed into it. Predictions for clustered search based on branch measurements
from a single flat pool must preserve this distinction.

For A, local binary search is exact. For unlimited B with valid bounds, it is exact as well.
For full-key C with complete weighted ordering, it is exact. Their global recall can still
be limited by routing. Finite node budgets and short-key quotas add a second possible loss.

\Needspace{12\baselineskip}
\subsection{Equivalent candidate sets}
The probability calculation describes a set of documents. Several algorithms can reach that
same set. If they score those documents with the same full score and use the same tie rule,
they have the same recall. Their time and memory costs can still be very different.

\textbf{A: open clusters and scan their documents.} For the random-bit example, use the first
$h$ bits as the cluster key. There are $C=2^h$ possible clusters. Probing all keys through
radius $t$ gives $P=H(h,t)$ clusters and exactly the candidate rule analyzed above. With
$h=4$ and $t=1$, scanning five clusters returns an expected 80 documents and has expected
top-one recall 89.60\%. The ratio $P/C=5/16$ describes the expected fraction of documents read;
it is not the nearest-document recall.

A k-means router uses distances to learned cluster centers instead of counting prefix errors. For
each true neighbor, let $J$ be the rank of its assigned cluster in the query's routing order.
Probing the first $P$ clusters retains that neighbor when $J\le P$. Routing recall is therefore
\[
r_{\rm route}(P)=\Pr(J\le P),
\]
where the probability averages over queries and a uniformly selected member of each query's
true top $K$. This identity defines what to measure. Cluster count and average cluster size
alone do not supply the distribution of $J$.

\textbf{B: follow bitplane branches.} Consider a four-level tree where a path specifies the
first four bits. Keep every path with at most one disagreement with the query. It reaches the
same five keys as the example above, so its recall is also 89.60\%. If every allowed path is
represented, the tree has $1+2+3+4+5=15$ nodes, including the root and five leaves. Some nodes
may have no documents and can be skipped.

This example counts a particular complete set of paths. The actual best-first algorithm has
a different stopping rule. It chooses a pending node using its score bound, splits on a query
coordinate, and may stop when the node budget is used. Its candidate set depends on the weights,
occupied leaves, leaf threshold, current result threshold, and visit order. A budget of 15
nodes does not mean that it reaches the five prefix keys above. To estimate its recall, run
that selection process on sampled data or derive its visited paths for a small explicit tree.

A random exploration probability changes which pending node is chosen. At a fixed budget it
can recover a missed neighbor or spend work on an unhelpful path. Its effect must be checked
with the same queries and budget. Complete traversal with valid bounds is the exact control.

\textbf{C: try keys near the query's desired pattern.} With equal weights, opening every short
key through a Hamming radius gives Equations~\eqref{eq:finite-top-one} and
\eqref{eq:finite-top-k}. With unequal weights, keys must be ordered by the sum of the weights
of their wrong bits. Several low-weight mismatches can precede a single high-weight mismatch.
The weighted distribution in the next subsections replaces the simple mismatch count.

For full keys, an exact weighted order and a valid stopping rule give exact local binary
results. Short keys omit part of the score, so stopping after enough postings have arrived can
lose neighbors. It can also achieve high recall when the prefix strongly predicts the full ranking.
The calculation must account for both possibilities.

Literal prefix relaxation is another concrete rule. For query key $0000$, replacing the last
constraint by a wildcard accepts both $0000$ and $0001$. A prefix trie can retrieve that parent
range, or only the new sibling range if the first key was already read. Flipping the last bit
instead visits only $0001$. Relaxing arbitrary query-selected coordinates is not a single
ordinary prefix-trie parent step. The candidate sets and data access costs must follow the
chosen structure.

\begin{examplebox}{Conditional acceptance across two checks}
Keep $d=16$ and suppose we know the nearest document has two errors. A router that demands the
first two bits match retains it with probability $\binom{14}{2}/\binom{16}{2}=91/120$.
Given that it passed, demanding the next two bits match has probability
$\binom{12}{2}/\binom{14}{2}=66/91$. Their product is $66/120=55\%$.

Multiplying $91/120$ by itself would give a different answer. Once the first two bits are known
to match, the two errors must lie among the remaining fourteen positions. Knowing the first check passed changes
the second calculation. The same issue arises when routing and lookup use overlapping bits.
\end{examplebox}

\Needspace{12\baselineskip}
\subsection{Omitted-score bounds}
Split the full score into retained and omitted coordinates:
\begin{equation}
S=T_h+Z_h,\qquad T_h=\sum_{i\in I_h}q_i b_i,\qquad
Z_h=\sum_{i\notin I_h}q_i b_i.
\end{equation}
Because $b_i\in\{-1,+1\}$,
\begin{equation}
|Z_h|\leq A_{\rm tail}=\sum_{i\notin I_h}|q_i|.
\label{eq:tail-error}
\end{equation}
Two documents can reverse their prefix order if the omitted contributions overcome their
prefix score difference. A prefix gap greater than $2A_{\rm tail}$ is sufficient to preserve
that pair's order. This is a sufficient bound, not a typical-error estimate.

A query-specific check compares the best unseen prefix score plus
$A_{\rm tail}$ with the worst fully scored retained result, as in Equation~\eqref{eq:prefix-bound}.
This condition depends on the actual query. If the bound is too loose, more keys must be visited. Returning
a fixed candidate count without checking that condition remains approximate.

When we omit coordinates from float vectors, the score error obeys the Cauchy--Schwarz inequality:
\[
|q_{\rm tail}^{\top}x_{\rm tail}|
\leq \|q_{\rm tail}\|_2\|x_{\rm tail}\|_2.
\]
Omitting coordinates and converting values to binary are separate approximations. Matryoshka training can improve
the usefulness of prefixes, but it does not set either bound to zero or supply a universal
recall percentage.

\Needspace{12\baselineskip}
\subsection{Weighted mismatches}
Let $E_i$ be one when document bit $i$ disagrees with the query and zero otherwise. With
$a_i=|q_i|$, the penalty is $L=\sum_i a_i E_i$. The score is $\sum_i a_i-2L$, so ranking by
largest score is the same as ranking by smallest $L$.

For query magnitudes $[4,3,2,1]$, one error in the first coordinate costs 4. Errors in the
last two cost $2+1=3$. The two-error document ranks higher. Counting wrong bits alone would
reverse that order.

Under independent mismatch probabilities $p_i$, a particular error pattern $e$ has probability
\begin{equation}
\Pr(E=e)=\prod_{i=1}^{d}p_i^{e_i}(1-p_i)^{1-e_i}.
\label{eq:error-pattern-probability}
\end{equation}
For a small dimension, list the $2^d$ patterns, calculate each full penalty and prefix penalty,
and add their probabilities into the corresponding score groups. That directly constructs the
joint distribution needed for recall: the probabilities of each full and prefix penalty pair. It also works with dependent bits if we supply their
actual joint pattern probabilities instead of the product above.

To handle a finite top $K$, group patterns by each distinct full penalty $\ell$. Let
$p_{\le\ell}=\Pr(L\le\ell)$ and $p_{<\ell}=\Pr(L<\ell)$. For a fixed candidate acceptance event
$A$, the same reasoning as the equal-weight case gives
\begin{equation}
\E[\operatorname{Recall@}K]
=\frac1K\sum_{\ell}
 \left[G_K(p_{\le\ell})-G_K(p_{<\ell})\right]\Pr(A\mid L=\ell).
\label{eq:weighted-finite-recall}
\end{equation}
Each bracket counts expected top-$K$ rows at that penalty. The conditional probability says
what fraction pass, given that full penalty. This is exact for independently sampled documents from the stated pattern
distribution, with an ID tie rule independent of the sampled patterns. It does not require
independence \emph{between bits}. It does require the joint distribution to be known.

\begin{examplebox}{Balanced bits can still be dependent}
Make eight-bit documents by sampling four balanced bits and repeating them in the last four
positions. Every individual bit is still zero or one half the time. However, only sixteen
of the possible 256 codes can occur. The independent-bit model assigns probability to 240
impossible codes. A correction to the arithmetic cannot repair that assumption; the pattern
probabilities must change. Section~\ref{sec:recall-checks} checks this example explicitly.
\end{examplebox}

At hundreds of dimensions, listing every pattern is impractical. We can sample from a stated
model, use a suitable dynamic program, or make an approximation to the score distribution.
The required calculation determines the appropriate method. The small exact model remains useful for
checking whether a faster calculation is answering the same question.

\Needspace{12\baselineskip}
\subsection{Independent-bit score correlation}
The following models initially exclude routing. When routing is added, retain the same global true-neighbor event or apply Equation~\eqref{eq:recall-stages}; replacing it with a different local top-$K$ event changes the question. When computing correlation, assume the scores have nonzero variance, meaning they do not all have the same value.

Hold one query $q$ fixed. Assume document signs are independent and equally
likely to be $-1$ or $+1$. Then
\begin{align}
\E[S\mid q]&=0,&\Var(S\mid q)&=\sum_{i=1}^{d}q_i^2,\\
\E[T_h\mid q]&=0,&\Var(T_h\mid q)&=\sum_{i\in I_h}q_i^2,\\
\Cov(T_h,S\mid q)&=\Var(T_h\mid q).&&
\end{align}
Correlation measures how strongly two scores vary together. Under these assumptions, it is exactly
\begin{equation}
\rho_q=\Corr(T_h,S\mid q)=
\sqrt{\frac{\sum_{i\in I_h}q_i^2}{\sum_{i=1}^{d}q_i^2}}.
\label{eq:rho-independent}
\end{equation}
Keeping much of the score variance does not mean the omitted score has a small worst-case bound. Many small omitted coordinates can have a modest sum of squares but a large sum of absolute values. A prefix can therefore give useful approximate results while still needing many checks to prove they are exact.

Retaining a fraction $v$ of score variance gives correlation
$\sqrt v$, not recall $v$ and not recall $\sqrt v$.

\begin{examplebox}{Prefix recall under two random-data distributions}
With equal coordinate variances, a 24-coordinate prefix of a 256-coordinate query retains
roughly $24/256$ of expected score variance. If instead the first 24 coordinates are constructed
to hold 80\% of expected score variance, their score varies much more closely with the
full score. The controlled measurements in Section~\ref{sec:evidence} test both cases. The
second is a deliberately favorable distribution, not the definition of Matryoshka embeddings.
\end{examplebox}

If query coordinates are random normals $q_i\sim\mathcal N(0,\sigma_i^2)$, we can choose each coordinate's variance $\sigma_i^2$. The ratio of \emph{expected} prefix and full variances is
$\sum_{i\in I_h}\sigma_i^2/\sum_i\sigma_i^2$. For a specific sampled query, use its actual
Equation~\eqref{eq:rho-independent}. The expectation of a ratio need not equal the ratio of
expectations.

\Needspace{12\baselineskip}
\subsection{Dependent document bits}
A covariance matrix records how coordinates vary together. Its diagonal entries are individual variances. An off-diagonal entry describes whether two coordinates tend to move together or in opposite directions. Independent, balanced signs have unit diagonal entries and zero off-diagonal entries.

For a general binary document vector with mean $\mu$ and covariance matrix $\Sigma$, let
$q_h$ be $q$ with omitted coordinates set to zero. Then
\begin{equation}
\rho_q=\frac{q_h^\top\Sigma q}
 {\sqrt{(q_h^\top\Sigma q_h)(q^\top\Sigma q)}}.
\label{eq:rho-covariance}
\end{equation}
This includes dependencies between coordinates. Restricting documents to a cluster or filter can
change both $\mu$ and $\Sigma$. How coordinates vary together across the whole database can differ from how they vary
inside the query's selected clusters.

A learned prefix may preserve ranking information through these dependencies even when its
share of the sum of squared coordinate values is small. Conversely, a prefix can have a large sum of squared values and
still fail to identify the extreme nearest neighbors. Estimating covariance also does not
prove that the rare highest scores follow a normal distribution.

\Needspace{12\baselineskip}
\subsection{Nonuniform mismatch probabilities}
Let $X_i=\ind[b_i\ne t_i]$. If the $X_i$ are independent with the same disagreement probability
$p$, and all relevant weights equal $a$, then a mismatch count over $h$ coordinates satisfies
\begin{equation}
M_h=\sum_{i=1}^{h}X_i\sim\operatorname{Binomial}(h,p),\qquad
\pi_h=aM_h.
\label{eq:binomial}
\end{equation}
A Hamming radius $r$ visits
\begin{equation}
V_h(r)=\sum_{u=0}^{r}\binom{h}{u}
\label{eq:hamming-volume}
\end{equation}
keys, and its expected candidate count is
\begin{equation}
N_*\,\Pr(M_h\leq r)
=N_*\sum_{u=0}^{r}\binom{h}{u}p^u(1-p)^{h-u}.
\label{eq:binomial-hits}
\end{equation}
Here $N_*$ is the relevant reference pool size; it is $N$ without filters or cluster restriction.
Equation~\eqref{eq:binomial-hits} is a candidate-count model, not yet recall.

If $p=1/2$, every key has equal probability and the expected count reduces to $N_*V_h(r)/2^h$.
Different $p_i$ with equal weights give a Poisson-binomial count: a sum of independent
zero-or-one outcomes with different probabilities. Unequal weights give a weighted Bernoulli sum,
which also multiplies each outcome by its weight. Neither is described by one ordinary binomial with an average $p$.

\Needspace{12\baselineskip}
\subsection{Recall under penalty thresholds}
Write $F_X(x)=\Pr(X\leq x)$ for a cumulative distribution function (CDF): the probability that $X$ is at or below a threshold.

Let $A_h$ be the prefix penalty and $B_h$ the omitted-coordinate penalty. Full penalty is
$A_h+B_h$. Under coordinate independence, the two sums are independent. Let $u$ be the
full-penalty threshold for a true neighbor, and let $v$ be the candidate prefix-penalty threshold. Then
\begin{equation}
r_{\rm population}(u,v)=
\frac{\Pr(A_h\leq v,\ A_h+B_h\leq u)}{\Pr(A_h+B_h\leq u)}
=\frac{\displaystyle\int_{a\leq v}F_{B_h}(u-a)\,dF_{A_h}(a)}
 {F_{A_h+B_h}(u)}.
\label{eq:population-recall}
\end{equation}
The numerator is a weighted sum. For each allowed prefix penalty, it multiplies the probability
that the remaining coordinates satisfy the full-score threshold by the probability of that
prefix penalty, then sums the contributions.

A candidate-count calculation alone omits the second condition: a retrieved document counts
toward recall only if its full penalty is also low enough to qualify.

For equal weights and an integer full radius $r_d$, the formula becomes
\begin{equation}
\frac{\displaystyle\sum_{u=0}^{\min(h,r_h,r_d)}
 \binom{h}{u}p^u(1-p)^{h-u}
 F_{\operatorname{Bin}(d-h,p)}(r_d-u)}
 {F_{\operatorname{Bin}(d,p)}(r_d)}.
\label{eq:binomial-recall}
\end{equation}
This is exact for the stated population events. Choosing those radii to approximate top $K$
and a candidate budget is another step. We must also account for ranks in a finite sample and score ties.

\begin{examplebox}{Candidate fraction is not recall}
Use all 64 six-bit codes as documents and an all-positive equal-weight query. The best seven
codes have zero or one mistake. Keep only codes matching the first two query bits: 16 codes
pass, so the candidate fraction is $16/64=25\%$. Five of the true seven pass: the ideal code
and the four codes whose one error is in an omitted coordinate. Recall is $5/7$, about 71.4\%,
not 25\%. Allowing one error in the two-bit prefix retrieves 48 codes and includes all seven.
Both numbers follow by counting the same true-neighbor event.
\end{examplebox}

\Needspace{12\baselineskip}
\subsection{Weighted key counts}
For a fixed query and lookup coordinates $I_h$, define
\begin{equation}
V_q(v)=\#\left\{S\subseteq I_h:\sum_{i\in S}|q_i|\leq v\right\}.
\label{eq:weighted-volume}
\end{equation}
This counts every attempted key through threshold $v$, whether occupied or empty. Under
independent uniform keys, expected candidates are $N_*V_q(v)/2^h$. With real key probabilities,
\begin{equation}
\E[F(v)\mid q]=N_*\sum_{k:\pi_q(k)\leq v}\Pr(\text{key}=k\mid q,\text{cluster},\text{filter}).
\label{eq:general-hits}
\end{equation}
For a clustered lookup, apply Equation~\eqref{eq:general-hits} separately with pool size $e_j$ and the key probabilities within cluster $j$, then sum the results.

An average bit-disagreement rate does not determine the probabilities of complete bit patterns.
The empty-key count $V_q$ and occupied candidate count $F$ are separate inputs to the latency
formula. A dense document distribution near the query key can make $F$ large while $V_q$
remains small; a sparse key space can do the reverse.

For small $h$, split weights in two, enumerate each half's subset sums, sort one list, and
count pairs whose sum is at most $v$. This meet-in-the-middle calculation uses roughly
$O(2^{h/2})$ memory and $O(h2^{h/2})$ sorting/search work. It is an analysis tool for counting
keys, not the measured latency of an online hash enumerator.

For weights $a_i$ rounded to nonnegative integers, a dynamic program builds the probability of each penalty from the previous step:
\begin{equation}
f_{i}(s)=(1-p_i)f_{i-1}(s)+p_i f_{i-1}(s-a_i).
\label{eq:dp}
\end{equation}
Replacing probabilities by counts and adding the two counts from the previous step computes the number of keys.
This is exact for the quantized-weight model. It is not an exact calculation
for unquantized query weights. If score quantization incurs a uniform score error at most
$\epsilon$, a score gap greater than $2\epsilon$ is sufficient to preserve a pair's order.
Otherwise candidates near the boundary need exact scoring or a safe expanded bound.

\Needspace{12\baselineskip}
\subsection{Normal score approximation}
When a score adds many small independent contributions, a normal distribution can sometimes
approximate its shape. Recall requires a joint distribution: how the full score and prefix score
vary together. This is an additional assumption, not a consequence of knowing their correlation. Subtract each score's mean and divide by its standard deviation. This puts them on the same scale as $(Z,Y)$ with correlation $\rho$. For a true
fraction $a=K/N_*$ and selected candidate fraction $f$, set
\[
z_K=\Phi^{-1}(1-a),\qquad y_f=\Phi^{-1}(1-f),
\]
where $\Phi$ and $\phi$ are the standard normal CDF and density. Under the joint-normal assumption, $Y$, when we know $Z=z$, has mean $\rho z$ and variance
$1-\rho^2$. Its probability of exceeding the prefix threshold is therefore
$\Phi((\rho z-y_f)/\sqrt{1-\rho^2})$. Multiply by the full-score density $\phi(z)$,
add across all scores that qualify as true neighbors, then divide by their probability $a$. The implied recall is
\begin{equation}
\widehat r=
\frac{1}{a}\int_{z_K}^{\infty}
\phi(z)\,\Phi\!\left(\frac{\rho z-y_f}{\sqrt{1-\rho^2}}\right)dz.
\label{eq:normal-recall}
\end{equation}
We average over \emph{all} scores that qualify as true neighbors. Substituting only $z=z_K$ calculates the passing probability at that one boundary,
not the average over all qualifying scores. When $\rho=0$, recall is $f$. In the limit of
identical rankings, $\rho=1$, it is $\min(1,f/a)$, subject to ties.

This formula is not valid simply because an embedding is called Matryoshka. A short weighted
binary prefix has only a finite set of possible scores. The probabilities of its highest scores can differ greatly
from those predicted by a normal distribution. At a billion documents, top-$K$ events may be much rarer than
anything covered by a small validation sample. Scores that vary closely together overall do not prove 99\% recall.

The controlled study below estimates score correlation on one query group and checks predicted
recall on different queries. Even there, the real-data normal approximation misses by several
percentage points. Such a model can suggest explanations for further testing; it cannot establish which method
is fastest at 99\% recall on real data.

\Needspace{12\baselineskip}
\subsection{Implications for the three methods}
\begin{itemize}
\item \textbf{A:} exact local scanning has no local candidate-search loss; global loss comes
from routing, candidate/reranker score changes, and the declared filtering task.
\item \textbf{B:} exact bound-based traversal has the same binary result as A on the same
selected pool. Finite budgets require query-level measurements or a distribution model for
visited nodes and surviving true neighbors. A node count alone is not enough.
\item \textbf{C:} full keys with complete weighted ordering are exact locally. Short keys can
be accurate when their ranking preserves true neighbors and the candidate budget is sufficient.
This possibility requires measurement. Their enumeration and posting costs remain
separate from the recall benefit.
\end{itemize}
A useful estimate needs a stated candidate rule and a distribution that fits the relevant
neighbors. The checks below show both a case where the exact calculation works and a case
where a plausible assumption produces an incorrect prediction.
